\section{An equivalent inverse using only rotations and XOR}
\label{sec:sm4_inverse_rotations}
The inverse matrix is one representation of $L^{-1}$, but it is not the only
way to evaluate it. For the SM4 data transformation, an equivalent formula is
\[
\boxed{L^{-1}(y)=\operatorname{ROR}_8\bigl(L(L(L(y)))\bigr).}
\]
Apply the normal, forward $L$ three times, using each result as the next input,
then rotate the final word right by eight bits. These are local word operations,
not three cipher rounds or three fault observations. On 32-bit words a right
rotation by 8 equals a left rotation by 24.

\subsection{Derivation from the SM4 definition}
Let $\rho_n$ denote a 32-bit left rotation by $n$, and let $I$ denote the identity
map. The symbols $\rho_n$ refer to rotations, not the output state words $R_j$.
The definition of the data map is
\[
L=I\xor\rho_2\xor\rho_{10}\xor\rho_{18}\xor\rho_{24}.
\]
Rotations commute and their amounts add under composition:
$\rho_a\rho_b=\rho_{a+b}$, with indices reduced modulo 32.
In the composition $L^2$, every cross term appears twice and cancels over GF(2).
For example, $\rho_2\rho_{10}\xor\rho_{10}\rho_2
=\rho_{12}\xor\rho_{12}=0$. Therefore
\[
\begin{aligned}
L^2&=I\xor\rho_4\xor\rho_{20}\xor\rho_{36}\xor\rho_{48}\\
&=I\xor\underbrace{\rho_4\xor\rho_4}_{0}\xor\rho_{20}\xor\rho_{16}\\
&=I\xor\rho_{20}\xor\rho_{16}.
\end{aligned}
\]
Square this expression once more:
\[
\begin{aligned}
L^4&=(L^2)^2=I\xor\rho_{40}\xor\rho_{32}\\
&=\underbrace{I\xor I}_{0}\xor\rho_8=\rho_8.
\end{aligned}
\]
If $y=L(x)$, then $L^3(y)=L^4(x)=\operatorname{ROL}_8(x)$.
A right rotation by 8 recovers $x$, proving the displayed inverse formula for
all 32-bit inputs. This identity is specific to this $L$; it must not be assumed
for arbitrary matrices or for the key-schedule transformation $L'$.

\subsection{The same numerical example as the matrix method}
The worked fault example gives
\[
y=\Delta X_{35}=\hex{734E3D3D}.
\]
The previous inverse-matrix calculation returned $\Delta B=\hex{4F000000}$.
The alternative follows this sequence:
\begin{center}
\begin{tabular}{ll}\toprule
Operation & Result (hexadecimal)\\\midrule
Input $y$ & \hex{734E3D3D}\\
First forward application $L(y)$ & \hex{4F04BF00}\\
Second application $L(L(y))$ & \hex{9DA47A90}\\
Third application $L(L(L(y)))$ & \hex{0000004F}\\
Final right rotation by 8 & \hex{4F000000}\\
Inverse-matrix method & \hex{4F000000}\\\bottomrule
\end{tabular}
\end{center}
Each forward application uses
$L(v)=v\xor\operatorname{ROL}_2(v)\xor\operatorname{ROL}_{10}(v)
\xor\operatorname{ROL}_{18}(v)\xor\operatorname{ROL}_{24}(v)$.
For example, the first application expands to
\[
\begin{aligned}
L(\hex{734E3D3D})={}&\hex{734E3D3D}\xor\hex{CD38F4F5}\xor\hex{38F4F5CD}\\
&\xor\hex{F4F5CD38}\xor\hex{3D734E3D}\\
={}&\hex{4F04BF00}.
\end{aligned}
\]
All rotations in one application use that application's input $v$, rather than
its running XOR accumulator. The next application starts only after that XOR
has been completed.

\begin{center}
\begin{tikzpicture}[node distance=4mm,box/.style={draw,rounded corners,align=center,fill=pale,text width=24mm,minimum height=12mm,font=\small}]
\node[box] (a) {\hex{734E3D3D}\\input};
\node[box,right=of a] (b) {\hex{4F04BF00}\\after $L$};
\node[box,right=of b] (c) {\hex{9DA47A90}\\after $L$};
\node[box,right=of c] (d) {\hex{0000004F}\\after $L$};
\node[box,right=of d] (e) {\hex{4F000000}\\right rotate 8};
\draw[-{Stealth}] (a)--(b);\draw[-{Stealth}] (b)--(c);
\draw[-{Stealth}] (c)--(d);\draw[-{Stealth}] (d)--(e);
\end{tikzpicture}
\end{center}

The final rotation moves the last byte to the first position:
\[
\hex{00 00 00 4F}\ \xrightarrow{\operatorname{ROR}_8}\ \hex{4F 00 00 00}.
\]
The forward check gives $L(\hex{4F000000})=\hex{734E3D3D}$.
Both methods thus recover the same S-box \emph{output difference}
$B_f\xor B_c$, not the individual clean output $B_c$.
The matrix method is a general inversion procedure; the rotation formula
exploits a special algebraic property of SM4. No measured performance comparison
between the two implementations is claimed here.
